\section{INSERT, UPDATE, DELETE}

Tiga perintah DML utama selain SELECT adalah \code{INSERT} (menambah data), \code{UPDATE} (mengubah data), dan \code{DELETE} (menghapus data). Ketiganya merupakan inti dari operasi \textbf{CRUD} (\textit{Create, Read, Update, Delete}) yang menjadi fondasi hampir semua aplikasi berbasis database. Setiap operasi harus dilakukan dengan hati-hati --- terutama UPDATE dan DELETE --- karena dapat mempengaruhi banyak baris sekaligus jika klausa WHERE tidak spesifik \cite{mysqlDocs}.

\code{INSERT INTO} menambahkan baris baru ke tabel. Terdapat dua sintaks: menyebutkan kolom secara eksplisit (lebih aman dan jelas) atau tanpa menyebutkan kolom (nilai harus sesuai urutan seluruh kolom). \code{UPDATE} mengubah nilai kolom pada baris yang memenuhi kondisi WHERE. \code{DELETE FROM} menghapus baris yang memenuhi kondisi WHERE. Tanpa WHERE, UPDATE mengubah \textbf{semua} baris dan DELETE menghapus \textbf{semua} baris --- ini sangat berbahaya \cite{mysqlGettingStarted}.

\begin{webcode}[language=SQL, caption={INSERT, UPDATE, dan DELETE}]
-- INSERT: menambah data baru
INSERT INTO produk (nama, harga, stok, kategori)
  VALUES ('Laptop ASUS', 8500000, 25, 'elektronik');

-- INSERT multiple rows
INSERT INTO produk (nama, harga, stok, kategori) VALUES
  ('Mouse Wireless', 150000, 100, 'elektronik'),
  ('Kaos Polos', 75000, 200, 'pakaian'),
  ('Kopi Arabika 250g', 45000, 50, 'makanan');

-- UPDATE: mengubah data (selalu gunakan WHERE!)
UPDATE produk SET harga = 8000000, stok = 20
  WHERE id = 1;

-- DELETE: menghapus data (selalu gunakan WHERE!)
DELETE FROM produk WHERE id = 4;

-- HATI-HATI: tanpa WHERE akan menghapus SEMUA data!
-- DELETE FROM produk;  -- JANGAN lakukan ini!
\end{webcode}

\begin{figure}[ht]
\centering
\begin{tikzpicture}[
  box/.style={draw, rounded corners, minimum width=2cm, minimum height=0.9cm, align=center, font=\small, fill=#1},
  arrow/.style={-Stealth, thick},
  node distance=1.2cm
]
  \node[box=green!20] (ins) {\textbf{INSERT}\\Tambah baris baru};
  \node[box=blue!15, right=1.5cm of ins] (sel) {\textbf{SELECT}\\Baca/ambil data};
  \node[box=orange!20, right=1.5cm of sel] (upd) {\textbf{UPDATE}\\Ubah data};
  \node[box=red!15, right=1.5cm of upd] (del) {\textbf{DELETE}\\Hapus data};

  \node[box=yellow!20, below=1.8cm of sel, minimum width=2.8cm] (tbl) {Tabel MySQL};
  \node[box=yellow!20, right=0pt of tbl, minimum width=2.8cm] (tbl2) {};

  \draw[arrow] (ins) -- (tbl);
  \draw[arrow] (tbl) -- (sel);
  \draw[arrow] (upd) -- (tbl2);
  \draw[arrow] (del) -- (tbl2);
\end{tikzpicture}
\caption{Empat operasi CRUD pada tabel MySQL}
\end{figure}

